\centerline{\bf \Huge Показательный листочек}

\begin{flushright}
        \scriptsize{Эпизод 4. Дилемма дикобразов.}
\end{flushright}

\problem{1}
Покажите, что при всех натуральных $a, n>1$ число $\varphi\left(a^n-1\right)$ делится на $n$.

\problem{2}
Найдите все тройки взаимно простых в совокупности натуральных чисел $(a, b, c)$ таких, что для любого натурального $n$ число $\left(a^n+b^n+c^n\right)^2$ делится на $a b+b c+c a$.

\problem{3}
Даны взаимно простые числа $a$ и $b$. Докажите, что для любого нечетного делителя $d$ числа $a^{2^n}+b^{2^n}$ число $d-1$ делится на $2^{n+1}$.

\problem{4}
Докажите, что для натурального числа $a$ найдется такое натуральное число $b>a$, что $1+2^b+3^b$ делится на $1+2^a+3^a$.

\problem{5}
Докажите, что для любого натурального $n>1$, свободного от квадратов, существуют такие простое число $p$ и целое число $m$, что $n$ делится на $p$, а $p^2+p m^p$ делится на $n$.

\problem{6}
Существует ли такое натуральное число $n$, что $6^n-1$ имеет ровно $n$ натуральных делителей?

\problem{7}
Для простых чисел $p \neq q$ нашлись натуральные числа $a$ и $b$ такие, что $a^p-b^p$ делится на $q$. Докажите, что $a-b$ делится на $q$ или $q-1$ делится на $p$.

\problem{8}
Найдите все натуральные $n$, для которых числа $n$ и $2^n-1$ имеют один и тот же набор простых делителей.

\problem{9}
Дано простое число $p$. Пусть $n$ - наименьшее натуральное число, большее 1 , такое, что $n^6-1$ делится на $p$. Докажите, что тогда одно из чисел $(n+1)^6-1$ или $(n+2)^6-1$ также делится на $p$.

\problem{10}
Докажите, что для простого $p$ уравнение $a^p+b^p=c^p$ не имеет решений в натуральных числах при $a, b, c \leqslant p$.

\problem{11}
Найдите все натуральные $k$ такие, что произведение первых $k$ нечетных простых чисел, уменьшенное на 1 , является точной степенью натурального числа (большей, чем первая).

\problem{12}
Дано нечетное натуральное число $m$. Докажите, что при некотором натуральном $n$ у числа $m^n+n^m$ не меньше 2025 различных простых делителей.
